\chapter{Frontend Overview}
\label{ch:frontend}

The Kiss frontend is a single-page application built with plain HTML, CSS, and JavaScript --- no build step, no transpiler, and no framework dependencies beyond Kiss's own component library.  This chapter covers the architecture, loading mechanism, and conventions that every Kiss frontend screen follows.

\section{Architecture}
\label{sec:fe-arch}

The frontend is served as static files.  There is no server-side rendering, no template engine, and no build pipeline.  You write HTML, CSS, and JavaScript directly, and the browser executes them as-is.

Navigation is handled by dynamically loading screen HTML and JavaScript into the page via \texttt{Utils.loadPage()}\index{Utils.loadPage}.  There is no client-side routing library; the framework manages screen transitions directly.

The frontend depends on:
\begin{itemize}
  \item \textbf{DOMUtils}\index{DOMUtils} --- Kiss's native DOM API wrapper, used for any direct DOM manipulation that falls outside the component layer.  jQuery is \emph{not} used.
  \item \textbf{AG-Grid Community} --- for data grids (Chapter~\ref{ch:grids}).
  \item \textbf{CKEditor} --- for rich text editing.
  \item \textbf{Kiss component library} --- custom HTML tags for forms and dialogs.
\end{itemize}

\section{Script Loading Order}
\label{sec:fe-loading}

The \texttt{index.html}\index{index.html} file is the entry point, but it does very little itself: it contains a small inline bootstrap kernel whose only job is to keep \texttt{index.html} fresh and then hand off to \texttt{kiss/bootstrap.js}\index{bootstrap.js}, which performs the actual ordered loading.  \texttt{bootstrap.js} defines the loader functions used throughout the framework:

\begin{description}[style=nextline]
  \item[\texttt{getScript(url)}] Loads a single script file.  Returns a Promise.
  \item[\texttt{getScripts(urls)}] Loads multiple scripts in parallel using \texttt{Promise.all()}.
  \item[\texttt{addStylesheet(url)}] Appends a stylesheet \texttt{<link>} to the page head.
\end{description}

\noindent
Every one of these helpers routes its URL through \texttt{window.cacheBust()} before fetching it, so every downloaded resource --- scripts and stylesheets alike --- is version-busted.  The mechanics of that busting (and of the kernel that precedes \texttt{bootstrap.js}) are covered in Chapter~\ref{ch:caching-security}; this section is concerned only with \emph{what loads, in what order, and why}.

\noindent
\texttt{bootstrap.js}'s \texttt{loadUtils()} function runs the following sequence:

\begin{enumerate}
  \item \textbf{Stylesheets} are added first and do not block anything that follows: \texttt{normalize.css}, \texttt{normalize2.css}, \texttt{kiss/starter-theme.css}, \texttt{kiss/Utils.css}, and the AG-Grid theme (\texttt{kiss/ag-grid-custom.css}).
  \item \textbf{Third-party libraries} load next, in parallel, and \emph{are} awaited: the AG-Grid data-grid engine and CKEditor.  They must finish before the Kiss wrapper scripts that reference their globals (\texttt{AGGrid.js}, \texttt{Editor.js}) load later in the sequence.
  \item \textbf{\texttt{SystemInfo.js}}\index{SystemInfo.js} loads next, by itself.  It carries the per-deployment configuration --- including \texttt{SystemInfo.stateStore}, needed by \texttt{AppState.init()} in the next step.  Once it has loaded, \texttt{bootstrap.js} mirrors the cache-control values from the \texttt{index.html} \texttt{<meta>} tags onto \texttt{SystemInfo.softwareVersion} / \texttt{SystemInfo.controlCache}, so existing references to those properties keep working.
  \item \textbf{\texttt{DOMUtils.js}}\index{DOMUtils.js} loads next, by itself.  Every other Kiss script that touches the DOM depends on it, and \texttt{Utils.js} in particular checks for the \texttt{DOMUtils} object at load time.
  \item \textbf{\texttt{AppState.js}}\index{AppState} loads next, by itself, and \texttt{AppState.init()} runs immediately afterward --- before \texttt{Server.js} loads, since \texttt{Server}'s \texttt{uuid} accessor reads through \texttt{AppState}.  Chapter~\ref{ch:routing-state} covers what \texttt{AppState} does.
  \item \textbf{The remaining framework scripts} load in parallel: \texttt{Utils.js}, \texttt{DateUtils.js}, \texttt{DateTimeUtils.js}, \texttt{TimeUtils.js}, \texttt{NumberUtils.js}, \texttt{Server.js}, \texttt{Router.js}, \texttt{AGGrid.js}, \texttt{Editor.js}, \texttt{MutableString.js}.
  \item \textbf{\texttt{routes.js}}\index{routes.js}, the application-owned route table, loads next, by itself --- it must be in place before \texttt{index.js} calls \texttt{Router.start()}.  See Section~\ref{sec:rs-routes-file}.
  \item \textbf{\texttt{index.js}} loads last, and is \emph{not} awaited --- nothing later in \texttt{loadUtils()} depends on it finishing.
\end{enumerate}

\begin{lstlisting}[style=javascript]
async function loadUtils() {
    // Stylesheets
    addStylesheet("normalize.css");
    addStylesheet("normalize2.css");
    addStylesheet("kiss/starter-theme.css");
    addStylesheet("kiss/Utils.css");
    addStylesheet("kiss/ag-grid-custom.css");

    // Third-party libraries -- must load before the kiss
    // AGGrid/Editor wrappers that use them
    await getScripts([
        "lib/ag-grid-community.min.noStyle.js",
        "lib/ckeditor.js"
    ]);

    // Per-deployment configuration (needed by AppState.init)
    await getScript("SystemInfo.js");
    SystemInfo.softwareVersion = window.KissVersion;
    SystemInfo.controlCache = window.KissCacheOn;

    // DOMUtils first (must complete before Utils.js); AppState
    // before Server.js uses it
    await getScript("kiss/DOMUtils.js");
    await getScript("kiss/AppState.js");
    AppState.init();

    await getScripts([
        "kiss/Utils.js",
        "kiss/DateUtils.js",
        "kiss/DateTimeUtils.js",
        "kiss/TimeUtils.js",
        "kiss/NumberUtils.js",
        "kiss/Server.js",
        "kiss/Router.js",
        "kiss/AGGrid.js",
        "kiss/Editor.js",
        "kiss/MutableString.js"
    ]);

    // routes.js (application-owned) registers routes and must
    // load before index.js calls Router.start()
    await getScript("routes.js");
    getScript("index.js");
}
\end{lstlisting}

\begin{warnbox}
If \texttt{DOMUtils.js} has not finished loading before \texttt{Utils.js} runs, you will see: ``DOMUtils object not found --- DOMUtils.js should be loaded before Utils.js.''  This is a race condition caused by loading them in parallel.  The sequence above avoids this, and the other ordering dependencies (third-party libraries before \texttt{AGGrid.js}/\texttt{Editor.js}, \texttt{SystemInfo.js} before \texttt{AppState.init()}, \texttt{AppState.js} before \texttt{Server.js}, \texttt{routes.js} before \texttt{index.js}), by \texttt{await}-ing each dependency before the scripts that need it are loaded.
\end{warnbox}

\section{Component Registration}
\label{sec:fe-components}

Before the login page can load, all UI components must be registered.  This happens in \texttt{index.js}:

\begin{lstlisting}[style=javascript]
(function () {
    Utils.useComponent('Popup');
    Utils.useComponent('CheckBox');
    Utils.useComponent('DateInput');
    Utils.useComponent('DropDown');
    Utils.useComponent('ListBox');
    Utils.useComponent('NumericInput');
    Utils.useComponent('PushButton');
    Utils.useComponent('RadioButton');
    Utils.useComponent('TextboxInput');
    Utils.useComponent('TextInput');
    Utils.useComponent('TextLabel');
    Utils.useComponent('TimeInput');
    Utils.useComponent('FileUpload');
    Utils.useComponent('NativeDateInput');
    Utils.useComponent('Picture');
})();
\end{lstlisting}

\noindent
Each \texttt{Utils.useComponent()} call loads a component's JavaScript from \texttt{kiss/component/<name>/}.  After registration, the corresponding custom HTML tags (e.g., \texttt{<text-input>}, \texttt{<push-button>}) become functional.

\section{The \texttt{\$\$()} Function}
\label{sec:fe-dollardollar}

The \texttt{\$\$()}\index{\$\$()} function is the primary way to access Kiss components.  Pass it the HTML \texttt{id} of a component and it returns the component's Kiss object with methods like \texttt{getValue()}, \texttt{setValue()}, \texttt{clear()}, \texttt{enable()}, \texttt{disable()}, and event handlers:

\begin{lstlisting}[style=javascript]
// Set a value
$$('employee-name').setValue("Alice");

// Get a value
let name = $$('employee-name').getValue();

// Clear a component
$$('employee-name').clear();

// Enable/disable
$$('save-btn').enable();
$$('save-btn').disable();

// Validate (shows error message if invalid, returns true
// on error)
if ($$('employee-name').isError('Employee Name'))
    return;

// Set focus
$$('employee-name').focus();

// Event handling
$$('save-btn').onclick(() => saveRecord());
$$('department').onChange(() => filterByDept());
\end{lstlisting}

\begin{warnbox}
Always use \texttt{\$\$()} and Kiss component methods for events --- not native DOM methods like \texttt{addEventListener()}.  The Kiss methods handle event wiring consistently across all components.
\end{warnbox}

\section{Screen Loading}
\label{sec:fe-screens}

Screens are pairs of \texttt{.html} and \texttt{.js} files in the \texttt{screens/} directory.  Each screen lives in its own subdirectory:

\begin{lstlisting}
src/main/frontend/screens/
  CRUD/
    CRUD.html
    CRUD.js
  Users/
    Users.html
    Users.js
  Framework/
    Framework.html
    Framework.js
\end{lstlisting}

\noindent
Load a screen with \texttt{Utils.loadPage()}:

\begin{lstlisting}[style=javascript]
Utils.loadPage("screens/CRUD/CRUD", "app-screen-area",
    "first-name");
\end{lstlisting}

\noindent
The three parameters are:
\begin{enumerate}
  \item \textbf{Path} --- the path to the screen files (without extension).  Kiss appends \texttt{.html} and \texttt{.js} automatically.
  \item \textbf{Target} --- the ID of the HTML element where the screen's HTML will be injected.  If omitted, the body is used.
  \item \textbf{Focus} --- the ID of a component that should receive focus after loading.
\end{enumerate}

\noindent
When a screen loads:
\begin{enumerate}
  \item The HTML file is fetched and injected into the target element.
  \item \texttt{Utils.rescan()} processes the injected HTML to activate any custom components.
  \item The JavaScript file is loaded and executed.
  \item Focus is set to the specified component.
\end{enumerate}

\section{Screen JavaScript Pattern}
\label{sec:fe-screen-js}

Every screen JavaScript file follows the same pattern --- an immediately-invoked async function expression (IIFE) that provides variable scope isolation:

\begin{lstlisting}[style=javascript]
'use strict';

(async function () {
    const WS = 'services/MyService';  // backend service path

    // Initialize grids, load data, set up event handlers
    const grid = new AGGrid('my-grid', columnDefs, 'id');
    grid.show();

    // Load initial data
    const res = await Server.call(WS, 'list');
    if (res._Success)
        grid.addRecords(res.records);

    // Button handlers
    $$('save-btn').onclick(async () => {
        // validate, call backend, refresh
    });
})();
\end{lstlisting}

\begin{tipbox}
The IIFE pattern ensures that each screen's variables are scoped to that screen.  When a new screen loads, the previous screen's variables are garbage-collected.
\end{tipbox}

\section{The Login Flow}
\label{sec:fe-login}

The application starts at \texttt{login.html}, which presents a login form:

\begin{lstlisting}[style=html]
<form onsubmit="return false;" autocomplete="on">
    <text-input id="username" required
                autocomplete="username"></text-input>
    <text-input id="password" required password
                autocomplete="current-password"></text-input>
    <push-button id="login">Login</push-button>
</form>
\end{lstlisting}

\noindent
The \texttt{<form>} tag with \texttt{onsubmit="return false;"} prevents actual form submission while enabling browser password managers via \texttt{autocomplete} attributes.

The \texttt{login.js} file:
\begin{enumerate}
  \item Validates the username and password fields.
  \item Calls \texttt{Server.call('', 'Login', data)} to authenticate.
  \item On success, stores the session UUID with \texttt{Server.setUUID(res.uuid)}.
  \item Enables back-button protection with \texttt{DOMUtils.preventNavigation()}.
  \item Loads the main application shell via \texttt{Utils.loadPage()}.
\end{enumerate}

\begin{lstlisting}[style=javascript]
const data = {
    username: $$('username').getValue().toLowerCase(),
    password: $$('password').getValue()
};
const res = await Server.call('', 'Login', data);
if (res._Success) {
    Server.setUUID(res.uuid);
    DOMUtils.preventNavigation(
        function () { return true; },
        function () {
            Utils.yesNo('Confirm',
                'Are you sure you want to logout?',
                function () { Server.logout(); });
        }
    );
    Utils.loadPage('screens/Framework/Framework');
}
\end{lstlisting}

\section{The Framework Screen}
\label{sec:fe-framework}

After login, the \texttt{Framework} screen provides the application shell: a navigation menu and a content area where individual screens are loaded.

\begin{lstlisting}[style=javascript]
// Framework.js - navigation button handlers
$$('crud-btn').onclick(function () {
    Utils.cleanup();
    Utils.loadPage('screens/CRUD/CRUD', 'app-screen-area');
});

$$('users-btn').onclick(function () {
    Utils.cleanup();
    Utils.loadPage('screens/Users/Users', 'app-screen-area');
});

$$('logout').onclick(function () {
    Server.logout();
});
\end{lstlisting}

\noindent
\texttt{Utils.cleanup()} removes event handlers and state from the current screen before loading a new one, preventing memory leaks.

\section{Backend URL Configuration}
\label{sec:fe-backend-url}

The \texttt{index.js} initialization code auto-detects the backend URL based on the deployment mode:

\begin{lstlisting}[style=javascript]
Utils.afterComponentsLoaded(function () {
    if (SystemInfo.backendUrl) {
        // explicit config
        Server.setURL(SystemInfo.backendUrl);
    } else if (SystemInfo.sameOriginBackend) {
        // served by the backend
        Server.setURL(Utils.getAppUrl());
    } else if (window.location.protocol === "file:") {
        // Electron desktop
        Server.setURL('http://localhost:8001');
    } else if (window.location.port >= 8000) {
        // development: backend is one port above
        Server.setURL('http://' +
            window.location.hostname + ':' +
            (Number(window.location.port) + 1));
    } else {
        // production
        Server.setURL(Utils.getAppUrl());
    }
});
\end{lstlisting}

The resolution order is:

\begin{enumerate}
  \item \textbf{Explicit configuration.}  \texttt{SystemInfo.backendUrl}, when set, always wins.  Use it for split deployments, or whenever the conventions below do not apply.
  \item \textbf{Served by the backend itself.}  \texttt{SystemInfo.sameOriginBackend}\index{sameOriginBackend} is \texttt{false} in the source tree, but the build system stamps it \texttt{true} into every copy of the front end that Tomcat serves --- the production WAR (during \texttt{./bld war}) and the development \texttt{tomcat/webapps/ROOT} (during \texttt{develop}/\texttt{start-backend}).  A stamped page calls the backend at its own origin, which makes browsing the Tomcat origin directly, and any production port, work with no configuration.  Never set this by hand.
  \item \textbf{Electron.}  Pages loaded from \texttt{file://} default to \texttt{http://localhost:8001}, the backend of the default port block; an application with a different \texttt{portBase} sets \texttt{SystemInfo.backendUrl} instead.
  \item \textbf{Development.}  An unstamped page on a local dev port came from the front-end static server, and under the Kiss port-block convention (Section~\ref{sec:build-ports}) the backend runs one port above it.  This derivation works for any \texttt{portBase} with no front-end changes.
  \item \textbf{Production fallback.}  Same origin as the page.
\end{enumerate}

\section{Mobile Support}
\label{sec:fe-mobile}

Kiss includes a separate \texttt{mobile/} directory with mobile-optimized login and screen templates.  The \texttt{index.js} initialization detects screen size and loads the appropriate login page:

\begin{lstlisting}[style=javascript]
const screenPixels = screen.height * screen.width;
if (screenPixels < 600000)
    Utils.loadPage("mobile/login");
else
    Utils.loadPage('login');
\end{lstlisting}

\noindent
Mobile screens use the same component library and \texttt{Server.call()} communication pattern as desktop screens, but with layouts optimized for smaller displays.

\section{Cache Control}
\label{sec:fe-cache}

In development, you may want to disable browser caching to ensure that file changes are picked up immediately.  The \texttt{index.html} file supports cache busting through the \texttt{SystemInfo.controlCache} flag:

\begin{lstlisting}[style=javascript]
const SystemInfo = {
    version: "2026.01.01.1",
    // append ?ver=VERSION to script URLs
    controlCache: true
};
\end{lstlisting}

\noindent
When \texttt{controlCache} is \texttt{true}, the \texttt{getScript()} function appends a version query parameter to each script URL, forcing the browser to re-fetch after version changes.
